\begin{vidu} %[2P4-7.5K][Loại 1]
Một electron sau khi được gia tốc nhờ hiệu điện thế \bth{U=300\ V}, chuyển động song song với một dây dẫn thẳng dài và cách dây dẫn một khoảng \bth{a=0,8\ cm}. Tìm lực lorenxo tác dụng lên electron nếu dòng điện chạy trong dây dẫn có cường độ \bth{I=10\ A}.
\choice
{\True $4,1.10^{-16}\ N$}
{$7,2.10^{-15}\ N$}
{$4,0.10^{-16}\ N$}
{$3,3.10^{-14}\ N$}
\loigiai{\begin{itemize}
\item Cảm ứng từ do dòng điện thẳng gây ra tại điểm đặt electron:
\begin{align*}
B=2.10^{-7} \dfrac{I}{a};\quad \vec{B}\perp \vec{v}\ct{(Vì} \vec{v} \ct{song song với dây)};\quad v=\sqrt{\dfrac{2eU}{m}}
\end{align*}
\item Lực Lorentz tác dụng lên electron: \bs{f}{|e|vB=4,1.10^{-16}N}.
\end{itemize}}
\haidong{4}
\end{vidu}
\begin{vidu} %[2P4-7.3K][Loại 1]
	Một proton và một hạt $\alpha$ không có vận tốc đầu, được tăng tốc bởi cùng một hiệu điện thế. Sau đó proton bay vào miền có từ trường đều theo hướng vuông góc với các đường sức. Khi đó quỹ đạo proton và $\alpha$ là các đường tròn có bán kính lần lượt là $R_P = 30$ cm và $R_\alpha$. Biết khối lượng $m_p = 1,672. 10^{-27}$ kg, $m_\alpha = 6,688. 10^{-27}$ kg, Bán kính quỹ đạo của hạt nhân helium $R_\alpha$ là
	\choice
	{$R_\alpha = 21,2$ cm}
	{\True $R_\alpha = 42,4$ cm}
	{$R_\alpha = 30$ cm}
	{$R_\alpha = 15$ cm}
	\loigiai{\begin{itemize}
			\item Động năng của hạt sau khi được tăng tốc bởi U là: $\dfrac{1}{2}m{{v}^{2}}=\left| q \right|U \Rightarrow v=\sqrt{\dfrac{2qU}{m}}$
			\item $\Rightarrow R=\dfrac{mv}{\left| q \right|B}=\dfrac{\sqrt{2mU}}{B\sqrt{\left| q \right|}}$
			\item Điện tích hạt nhân helium là $q_\alpha = 2e$, điện tích hạt proton là $q_p = e.
			\Rightarrow \dfrac{{{R}_{\alpha}}}{R_p}=\sqrt{\dfrac{{{m}_{\alpha}}}{m_p}. \dfrac{q_p}{{{q}_{\alpha}}}}=\sqrt{2}\Rightarrow {{R}_{\alpha}}=30\sqrt{2}$ cm.
	\end{itemize}}
\haidong{5}
\end{vidu}
\begin{vidu} %[2P4-7.5G][Loại 2]
	Một proton chuyển động thẳng đều trong miền có cả từ trường đều và điện trường đều. Véctơ vận tốc của hạt hướng từ trái sang phải và đường sức điện trường hướng từ trong ra ngoài mặt phẳng giấy. Biết $E = 6000$ V/m, $v = 3.10^6$ m/s, xác định hướng và độ lớn $\vec{B}$ .
	\choice
	{$\vec{B}$ hướng lên. $B = 0,003$ T}
	{\True $\vec{B}$ hướng xuống. $B = 0,002$ T}
	{$\vec{B}$ hướng lên. $B = 0,004$ T}
	{$\vec{B}$ hướng xuống. $B = 0,0024$ T}
	\loigiai{Proton mang điện dương nên E cùng chiều với f. $\Rightarrow$ Áp dụng quy tắc bàn tay trái ta xác định được B hướng từ dưới lên.$f = |q|vB\sin \alpha = |q|E \Rightarrow B = E/(v.sin90) = 0,002$ T}
\haidong{3}
\end{vidu}
\begin{vidu} %[2P4-7.5K][Loại 3]
	Thành phần nằm ngang của từ trường trái đất bằng $3.{10^{- 5}}$ T, thành phần thẳng đứng rất nhỏ. Một proton chuyển động theo phương ngang theo chiều từ Tây sang Đông thì lực Lorenxơ tác dụng lên nó bằng trọng lượng của nó, biết khối lượng của proton là $1,67.{10^{- 27}}$ kg và điện tích là $1,6.{10^{- 19}}$ C. Lấy $g = 10\ m/s^2$, tính vận tốc của proton
	\choice
	{$3.{3.10^{- 3}}$ m/s}
	{$1,{5.10^{- 3}}$ m/s}
	{\True $3,{5.10^{- 3}}$ m/s}
	{$2,{5.10^{- 3}}$ m/s}
	\loigiai{\begin{itemize}
			\item Ta có $f = P = mg = 1,67.10-26{10^{- 26}}$
			\item Mà $v = \dfrac{f}{\left| q \right|B\sin \theta} = 3,5.10-3\ m/s.3,{5.10^{- 3}}\ m/s$.
	\end{itemize}}
\haidong{3}
\end{vidu}
\begin{vidu} %[2P4-7.5T][Loại 4]
	Một êlectrôn không vận tốc đầu sau khi được tăng tốc bởi hiệu điện thế tăng $U = 40$ V nó bay vào một vùng từ trường đều có hai mặt biên phẳng song song, bề dày $h = 10$ cm. Vận tốc của êlectrôn vuông góc với cả vector cảm ứng từ lẫn hai biên của vùng. Với giá trị nhỏ nhất $B_{\min}$ của cảm ứng từ bằng bao nhiêu thì êlectrôn không thể bay xuyên qua vùng đó? Cho biết tỷ số độ lớn điện tích và khối lượng của êlectrôn là $\mu = 1,76. 10^11$ C/kg.
	\choice
	{$3,2. 10^{-4}$ T}
	{$4,2. 10^{-4}$ T}
	{$2,8. 10^{-4}$ T}
	{\True $2,1. 10^{-4}$ T}
	\loigiai{\begin{itemize}
			\item Thế năng êlectrôn nhận được khi đi qua hiệu điện thế tăng tốc chuyển thành động năng của êlectrôn $eU=\dfrac{1}{2}m{{v}^{2}}$ $\Rightarrow v=2eUm=2\mu U \Rightarrow v=\sqrt{\dfrac{2eU}{m}}=\sqrt{2\gamma U}$
			\item Khi êlectrôn chuyển động vào vùng từ trường đều với vận tốc vuông góc với vector cảm ứng từ thì quỹ đạo chuyển động của êlectrôn là đường tròn bán kính $R=\dfrac{mv}{eB}=\dfrac{v}{\gamma B}$
			\item Để êlectrôn không thể bay xuyên qua vùng từ trường đó thì bán kính quỹ đạo R $\le$ h $\Rightarrow {{R}_{\max}}=\dfrac{v}{\gamma {{B}_{\min}}}=h\Rightarrow {{B}_{\min}}=\dfrac{1}{h}\sqrt{\dfrac{2U}{\gamma}}=2,{{1. 10}^{-4}}$ T
	\end{itemize}}
\haidong{8}
\end{vidu}

